% Inactive research fragment. It is not an input of the checked book.
\subsection{Curve backgrounds and finite curve factors}
\label{curve-backgrounds647:scope}

Let $p:X\to C$ be a surjective holomorphic map with connected
fibers, where $X$ is a connected compact complex manifold of
dimension $n\geq2$ and $C$ is a connected compact Riemann surface.
Let $\sigma$ be any smooth real $(1,1)$-form on $C$, and put
$A=\int_C\sigma$ and $\theta=p^*\sigma$. Critical fibers are
allowed. Fix a smooth positive area-one form $\eta$ on $C$.

\begin{proposition}\label{curve-backgrounds647:envelopes-and-mass}
The class $\PSH(X,\theta)$ is nonempty if and only if $A\geq0$.
For $v\in\PSH(C,\sigma)$ and $U=v\circ p$, one has
\[
 P_\theta[U]=P_\sigma[v]\circ p,\qquad
 P_\theta[U]_{\mathcal I}=P_\sigma[v]_{\mathcal I}\circ p.
\]
Pullback is a bijection between the ordinary model spaces and
between their ideal subspaces. Every potential in this background
has zero extended Monge--Amp\`ere measure, and every ordinary type
has zero upper and lower singularity volumes.

If $A=0$ and $\sigma=\ddc h$, the only normalized potential is
$-h\circ p+\min_Ch$. It is smooth and is the unique ordinary and
ideal model.
\end{proposition}
\begin{proof}
By \cref{curve-map:potential-descent}, every potential descends,
even for this possibly indefinite background. Thus nonemptiness
is equivalent to the curve condition in
\cref{lem:curve-background-area}. Descent preserves the ordinary
order, and \cref{curve-map:ideal-types} preserves and reflects
the all-weight ideal order. The map $p$ is open, as proved in
\cref{fibration:curve-envelope-pullback}. Hence upper
regularization commutes with pullback, proving the envelope
identities and the model bijections.

Write $\sigma=A\eta+\ddc h$ with $h$ smooth, using
\cref{lem:curve-background-area}. For $v\in\PSH(C,\sigma)$, put
$w=v+h\in\PSH(C,A\eta)$. Then
\[
 \theta+\ddc(v\circ p)
   =A p^*\eta+\ddc(w\circ p).
\]
When $A>0$, \cref{fibration:curve-envelope-pullback} applied to
the positive curve form $A\eta$ gives zero extended measure on
the right. The current invariance in
\cref{thm:arbitrary-extended-ma} gives zero extended measure in
the original decomposition as well. The smooth correction
reduces the calculation to admissible truncations in the
positive background.

When $A=0$, the curve lemma gives $v=-h+c$. Its pullback is
smooth with zero curvature current, so its top product is zero.
Surjectivity gives the stated normalization. Every potential
differs from this one by a constant, and both envelopes fix its
normalized representative. Since descent applies to every
representative of every ordinary type, both singularity volumes
are zero in all cases.
\end{proof}

Assume now that $A\geq0$. For $v\in\PSH(C,\sigma)$ write
\[
 T_v=\sigma+\ddc v,\qquad
 \mu_v=\mathbf1_{\{v=-\infty\}}T_v.
\]
Fix a Gauduchon form $\gamma$ on $X$. Let $L_\gamma$ and
$\kappa_\gamma$ be the pullback trace operator and the fiber
volume of \cref{fibration:gauduchon-fiber-volume}. These are
defined using $p$ and finite measures on $C$; they do not require
positivity of $\sigma$. For an ordinary model $U=v\circ p$, put
\[
 \lambda_U^\gamma
   =\mathbf1_{\{U=-\infty\}}
       (\theta+\ddc U)\wedge\gamma^{n-1},\qquad
 \rho_\theta(U,V)
   =\kappa_\gamma^{-1}
       \|\lambda_U^\gamma-\lambda_V^\gamma\|_{\mathrm{TV}}.
\]

\begin{theorem}\label{curve-backgrounds647:full-model-metric}
Pullback is an isometry from
$(\mathcal M(\sigma),d_\sigma)$ onto
$(\mathcal M(\theta),\rho_\theta)$. The latter space is complete
and geodesic, with diameter $2A$, and its metric is independent
of $\gamma$. The ideal models are a closed geodesically convex
subspace equal to the metric closure of the analytic models.
Ideal projection is its unique nearest-point retraction and is
$1$-Lipschitz. The distance of $v\circ p$ to this subspace is the
nonatomic polar mass of $v$.

Metric convergence implies $L^1(X)$ convergence for every smooth
positive volume form. The upper-volume comparison is identically
zero. At $A=0$ this is the one-point metric space.
\end{theorem}
\begin{proof}
The definition of pullback measure gives
$p^*T_v=\theta+\ddc(v\circ p)$: in the smooth decomposition
$\sigma=A\eta+\ddc h$, both sides are
$A p^*\eta+\ddc((v+h)\circ p)$. Apply the Borel restriction and
total-variation identities of
\cref{fibration:gauduchon-fiber-volume} to get
\[
 \lambda_{v\circ p}^\gamma=L_\gamma(\mu_v),\qquad
 \rho_\theta(v\circ p,w\circ p)
       =\|\mu_v-\mu_w\|_{\mathrm{TV}}=d_\sigma(v,w).
\]
The envelope bijection and
\cref{prop:curve-arbitrary-background-models} give the isometry,
completeness and geodesics. Polar measures have mass at most
$A$, so their mutual distance is at most $2A$. For $A>0$, the
models of $A\delta_s$ and $A\delta_t$, $s\ne t$, attain this
bound. For $A=0$, the preceding proposition gives the singleton.

Ideal projection corresponds to the atomic part of the polar
measure by the same curve classification. If
$\mu=\mu_{\mathrm{at}}+\mu_{\mathrm{diff}}$ and $\nu$ is atomic,
the measures $\mu_{\mathrm{diff}}$ and
$\mu_{\mathrm{at}}-\nu$ are mutually singular. Therefore
\[
 \|\mu-\nu\|_{\mathrm{TV}}
  =\mu_{\mathrm{diff}}(C)
     +\|\mu_{\mathrm{at}}-\nu\|_{\mathrm{TV}}.
\]
The minimum is attained exactly at $\nu=\mu_{\mathrm{at}}$.
Taking the atomic part does not increase total variation, giving
the Lipschitz estimate. Atomic polar measures of mass at most
$A$ form a closed convex subset of the polar-measure space.

For $A>0$, the smooth-gauge model bijection in the proof of
\cref{prop:curve-arbitrary-background-models} preserves analytic
singularities in both directions: its envelope differs by a
bounded function from the smoothly transported analytic model,
by \cref{prop:i-envelope-gauge}. Thus the analytic density in
\cref{cor:curve-model-completeness} carries from $A\eta$ to
$\sigma$. The analytic descent equivalence of
\cref{curve-map:analytic-types} and the envelope bijection then
give the stated closure upstairs. At zero area the unique smooth
model is already analytic and ideal.

Curve metric convergence implies $L^1(C)$ convergence by
\cref{prop:curve-arbitrary-background-models}, and
\cref{curve-map:lone-pullback} gives convergence upstairs.
Finally all top-degree masses, including those of pairwise
maxima, vanish by the preceding proposition. The upper-volume
comparison therefore vanishes.
\end{proof}

\begin{corollary}\label{curve-backgrounds647:all-obstacles}
For $U=v\circ p$ and $R=P_\theta[U]_{\mathcal I}$, the following
are equivalent: $P_\theta[U]=P_\theta[U]_{\mathcal I}$;
$\mu_v$ is atomic; $U\in\mathcal E_{\mathrm{env}}(X,\theta,R)$;
and the ordinary and ideal envelopes agree below every smooth
real obstacle on $X$. At zero area they all hold.
\end{corollary}
\begin{proof}
Put $r=P_\sigma[v]_{\mathcal I}$. The envelope bijection gives
$R=r\circ p$. By \cref{cor:curve-arbitrary-background-criterion},
base envelope equality is equivalent to atomic polar curvature
and to $v\in\mathcal E_{\mathrm{env}}(C,\sigma,r)$.
For each stretching parameter $b\geq1$, descent and surjectivity
identify competitors satisfying
$b^{-1}h+(1-b^{-1})r\leq v$ with their upstairs counterparts.
This uses only convex combinations and remains valid at poles.
Thus the relative-class conditions correspond. Apply
\cref{prop:ideal-obstacle-reprojection} for equality below every
smooth obstacle on $X$; the converse includes the zero obstacle.
The zero-area assertion follows from the smooth singleton.
\end{proof}

Keep $X$ and $C$ as above. We now allow $p:X\to C$ to have
disconnected fibers and suppose that it has a factorization
\[
 X\xrightarrow{q}C'\xrightarrow{r}C,
 \qquad p=r\circ q,
\]
where $q$ is surjective and holomorphic with connected fibers,
$C'$ is a connected compact Riemann surface, and $r$ is finite
and holomorphic of degree $d\geq1$. The factorization is part
of the hypotheses.
Let $\beta$ be any smooth real form on $C$ with
$A=\int_C\beta\geq0$, and put $\sigma=r^*\beta$ and
$\theta=p^*\beta=q^*\sigma$.

\begin{proposition}\label{curve-backgrounds647:finite-factor}
The entire ordinary $\theta$-model space is isometric, using the
$q$-fiber normalization, to the curve model space in background
$\sigma$, whose area is $dA$. Its diameter is $2dA$; it has the
analytic closure, ideal retraction, convergence and all-obstacle
properties above.

If $\kappa_q^\gamma$ and $\kappa_p^\gamma$ are the constants
defined by the pushforwards of $\gamma^{n-1}$ to $C'$ and $C$,
then
\[
 \kappa_p^\gamma=d\kappa_q^\gamma.
\]
Consequently normalizing the same polar-trace variation by
$\kappa_p^\gamma$ gives $d^{-1}$ times the $q$-normalized metric,
with diameter $2A$. This rescaled metric is also independent of
$\gamma$. For a model $U=v\circ q$ with polar measure $\mu_v$,
\[
 p_*\lambda_U^\gamma=\kappa_q^\gamma r_*\mu_v.
\]
\end{proposition}
\begin{proof}
Away from the finitely many critical values, $r$ is a covering
with $d$ sheets. Change of variables on these sheets gives
\[
 \int_{C'}r^*\alpha=d\int_C\alpha
\]
for every smooth real $(1,1)$-form $\alpha$; the branch points
and their inverse images have smooth area zero. In particular
$\int_{C'}\sigma=dA$. Apply the three preceding results to $q$
and this smooth curve background, which may vanish at branch
points or be indefinite.

The Gauduchon trace lemma for $q$ gives
$q_*(\gamma^{n-1})=\kappa_q^\gamma$. For a smooth test form
$\alpha$ on $C$, the projection formula and the covering
calculation yield
\[
 \int_Xp^*\alpha\wedge\gamma^{n-1}
  =\kappa_q^\gamma\int_{C'}r^*\alpha
  =d\kappa_q^\gamma\int_C\alpha.
\]
Thus $p_*(\gamma^{n-1})$ is the constant $d\kappa_q^\gamma$.
The ratio of the two normalizations is $d^{-1}$, giving every
metric assertion. Finally
$q_*\lambda_U^\gamma=\kappa_q^\gamma\mu_v$, so composing
pushforwards proves the last identity.
\end{proof}

\begin{example}\label{curve-backgrounds647:disconnected-fibers}
Let $X=\mathbb P^1\times\mathbb P^1$, let $q$ be its first
projection and let $r:\mathbb P^1\to\mathbb P^1$ be
$r([z_0:z_1])=[z_0^2:z_1^2]$. Put $p=r\circ q$ and
$\theta=p^*\omega_{\FS}$, with $\int_{\mathbb P^1}\omega_{\FS}=1$.
There are two distinct analytic ideal $\theta$-models $U_+,U_-$
which do not descend through $p$. Their $q$-normalized polar
distance is four and their $p$-normalized distance is two, yet
\[
 p_*\lambda_{U_+}^\gamma=p_*\lambda_{U_-}^\gamma
\]
for every Gauduchon form $\gamma$. Thus the pushforward of polar
trace to the original base does not separate the models.
\end{example}
\begin{proof}
The finite map $r$ has degree two and is branched at zero and
infinity. In the affine coordinate it is $r(z)=z^2$, so the
smooth semipositive form $\sigma=r^*\omega_{\FS}$ has area two.
For $s\in\{1,-1\}$, set
\[
 v_s(z)=2\log|z-s|^2-\log(1+|z|^4)-\log8.
\]
This extends smoothly across infinity and has curvature
$\sigma+\ddc v_s=2\delta_s$. The inequality
$|z-s|^4\leq(|z|+1)^4\leq8(1+|z|^4)$, with equality at
$z=-s$, shows that its supremum is zero.
The model of $2\delta_s$ from
\cref{prop:curve-arbitrary-background-models} has precisely
this total curvature, since its polar measure already has the
full mass two. Harmonic uniqueness and normalization identify
it with $v_s$. Write $v_+=v_1$ and $v_-=v_{-1}$. Both have
analytic singularities, hence are ideal by
\cref{prop:analytic-envelope-type}. The model pullback bijection makes
$U_\pm=v_\pm\circ q$ analytic ideal $\theta$-models.

The function $U_+$ is $-\infty$ on $q^{-1}(1)$ and finite on
$q^{-1}(-1)$, although both divisors are mapped by $p$ to the
same base point $1$. It therefore cannot descend through $p$;
the same argument with the divisors exchanged applies to $U_-$.
The variation of $2\delta_1-2\delta_{-1}$ is four. The finite
factor proposition gives the two distances and
\[
 p_*\lambda_{U_\pm}^\gamma
   =2\kappa_q^\gamma\delta_1.
\]
The upstairs polar traces are distinct, while this common
pushforward loses which component carries the pole.
\end{proof}
